\question[4]
Die folgenden beiden Zeilen zeigen den Ausgangspunkt für 
eine Vigenère-Verschlüsselung. Erläutere an dem Beispiel, wie die Verschlüsselung 
funktioniert, bestimme das Schlüsselwort und den resultierenden Geheimtext.


\texttt{INFORMATIK}  (Klartext)\\
\texttt{CODECODECO}  (Schlüssel)\\


\begin{solutionbox}{3cm}
 
Das Schlüsselwort ist \texttt{CODE}. Der Schüssel ist die Wiederholung des Schlüsselworts bis die
Länge des Klartexts erreicht ist. Der Klartext wird 
zeichenweise mit dem Schlüssel \texttt{CODECODECO} 
verschlüsselt. Dabei wird jeder Buchstabe des Klartexts um die Position des 
entsprechenden Buchstabens im Schlüssel verschoben.
Der resultierende Geheimtext ist \texttt{KBISTADXKY}.
 
\end{solutionbox}